% Exact image description; no surjectivity assertion is made here.
\subsection{The exact finite images of the descent endomorphism}

Write $B=\widehat{C_2(s)*C_3(r)}$, $\tau=sr$,
$F=\widehat F(x,y)\triangleleft B$, $x=\tau^{-2}$,
$y=sxs$, and let
\[
 H=\ker\bigl(\chi:\widehat{GT}_0\longrightarrow
                         \widehat{\mathbb Z}^{\times}\bigr).
\]
The quotient used throughout this subsection is the literal
map $\varphi_P:F\twoheadrightarrow B$, with
$\varphi_P(x)=s$, $\varphi_P(y)=\tau$.  All derived and
commutator subgroups in this subsection are closed.

We first specify finite sets and maps for arbitrary finite
levels.  Let $q_R:B\twoheadrightarrow R$ be a finite marked
quotient with kernel $K_R\subseteq F\cap B'$.  Its canonical
marked central cover and its pure subgroup are
\[
\begin{aligned}
 R^{\mathrm c}&=B/\overline{[B,K_R]},&
 M_R&=K_R/\overline{[B,K_R]},\\
 P_R&=q_R(F),&
 P_R^{\mathrm c}&=\operatorname{im}(F\longrightarrow R^{\mathrm c}).
\end{aligned}                                                   \tag{B.Q1}
\]
Here $R^{\mathrm c}$ is finite.  Indeed $M_R$ is a
topologically finitely generated central abelian group,
because $K_R$ is open in the finitely generated profinite
group $B$.  The containment $K_R\subseteq B'$ gives
$M_R\subseteq(R^{\mathrm c})'$.  Transfer to the central
subgroup $M_R$ of index $|R|$ kills $(R^{\mathrm c})'$ and
sends $m\in M_R$ to $m^{|R|}$.  Thus $M_R$ has exponent
dividing $|R|$ and is finite.  This is the same actual-source
argument as in B.2.6; it requires no solvability assumption
on $R$.

Denote the distinguished generator images in (B.Q1) by the
same lower-case letters as their source generators.  Define
the following finite group of framed automorphisms:
\[
\begin{split}
 \mathcal A(R^{\mathrm c})=\bigl\{(\alpha,\beta,A):\;&
 \alpha\in\operatorname{Aut}(R^{\mathrm c}),\quad
 \alpha(M_R)=M_R,\quad\alpha(\tau)=\tau,\\
 &\beta,A\in P_R^{\mathrm c},\quad
 \alpha(s)=\beta s\beta^{-1},\quad
 \alpha(r)=ArA^{-1}\bigr\}.
\end{split}                                                    \tag{B.Q2}
\]
Every such $\alpha$ acts trivially on
$R^{\mathrm c}/P_R^{\mathrm c}=S_3$, so it preserves
$P_R^{\mathrm c}$.  With automorphisms composed as functions,
the multiplication in (B.Q2) is
\[
 (\alpha,\beta,A)(\alpha',\beta',A')
   =(\alpha\alpha',\,\alpha(\beta')\beta,
                        \,\alpha(A')A).                       \tag{B.Q3}
\]
Define $\mathcal A(R)$ by the same equations in $R$, omitting
the condition on $M_R$.

The finite-certificate image equality B.2.7 applies verbatim
to (B.Q1): the image in $\mathcal A(R)$ of full continuous
positive automorphisms of $B$ fixing $\tau$ and preserving
$K_R$, with specified pure frames, equals the reduction image
of $\mathcal A(R^{\mathrm c})$.  To make the scope explicit,
the finite extension used in its proof is every joint
quotient $E\twoheadrightarrow R$, with the normal kernel
$\ker(E\to R)$.  The construction R1, (D.50.5)--(D.50.16),
uses its subgroup lattice and finite complex irreducible
representations without requiring solvability.  Its count
argument uses only normality, as the paragraph preceding
(D.50.17) explicitly states.  Every resulting marked
central quotient $R^\dagger\to R$ factors through
$R^{\mathrm c}$.  For $v:R^{\mathrm c}\twoheadrightarrow
R^\dagger$, the actual new source map is
$v\alpha q_{R^{\mathrm c}}$, an epimorphism with the three
positive frame values $v(\beta),v(A),1$.  The equality of
epimorphism counts supplies all finite conditions.  Compactness,
followed by the inner correction using the retained cusp
conjugator, gives the full $\tau$-fixed lift with its specified
$R$-frames.  This is precisely the proof of (B.M12), with $R$
in place of $D$; the separate joint-generation assertion
(D.50.17) is not used.

Let $\mathcal Q$ be the directed family of all quotients
\[
 q_Q:B\twoheadrightarrow Q=B/K_Q,
 \qquad K_Q\text{ open characteristic in }B,
 \qquad K_Q\subseteq F\cap B'.                                \tag{B.Q4}
\]
The ordering is by refinement.  This family is cofinal among
finite quotients of $B$: a finitely generated profinite group
has only finitely many open subgroups of each given index,
so every open subgroup contains an open characteristic
subgroup.  Apply this also to its intersection with $F\cap B'$.
In particular each quotient retains the specified maps to
$S_3=B/F$ and $C_6=B/B'$.

For each $Q\in\mathcal Q$, construct the following actual
finite source quotient and its actual pure quotient map:
\[
\begin{aligned}
 L_Q&=\ker(q_Q\varphi_P:F\longrightarrow Q),\\
 K_{E_Q}&=B'\cap\operatorname{core}_B(L_Q),
       &E_Q&=B/K_{E_Q},\\
 \Theta_Q:P_{E_Q}=F/K_{E_Q}&\twoheadrightarrow Q,
       &\Theta_Q(x)&=q_Q(s),\quad
         \Theta_Q(y)=q_Q(\tau).
\end{aligned}                                                   \tag{B.Q5}
\]
The core is the intersection of the six conjugates of $L_Q$
indexed by $B/F$.  Thus (B.Q5) is a finite quotient of the
specified source, obtained by retaining the six conjugate
maps $F\to Q$, their permutation action by $B/F$, and the
abelian map $B\to C_6$.  Its kernel is contained in
$F\cap B'$, so (B.Q1)--(B.Q3) apply to $E_Q$.

The normalized faithful representation
$H\to\operatorname{Out}(B)$ gives continuous finite maps
\[
 \rho_Q:H\longrightarrow\operatorname{Out}(Q).
\]
There are also finite maps
\[
\begin{aligned}
 p_Q:\mathcal A(Q^{\mathrm c})&\longrightarrow
                         \operatorname{Out}(Q),\\
 r_Q:\mathcal A(E_Q^{\mathrm c})&\longrightarrow
                         \operatorname{Out}(Q).
\end{aligned}                                                   \tag{B.Q6}
\]
The first is reduction of the automorphism coordinate followed
by passage to its outer class.  The second is reduction to
$E_Q$, restriction to $P_{E_Q}$, descent through $\Theta_Q$,
and passage to the outer class.  The proof below verifies that
this descent exists for every member of its stated domain;
no kernel-preservation hypothesis is omitted from (B.Q6).

\begin{result}{B.2.9}{Proved}{Exact finite descriptions of the descent image}
For every $Q\in\mathcal Q$, the maps (B.Q6) are well-defined
homomorphisms of finite groups, and their images are exactly
\[
\begin{aligned}
 \rho_Q(H)&=p_Q\bigl(\mathcal A(Q^{\mathrm c})\bigr),\\
 \rho_Q\bigl(D_P(H)\bigr)
          &=r_Q\bigl(\mathcal A(E_Q^{\mathrm c})\bigr).
\end{aligned}                                                   \tag{B.Q7}
\]
Consequently the full multiplier-one descent image is the
following closed subgroup of $H$:
\[
 D_P(H)=\bigcap_{Q\in\mathcal Q}
       \rho_Q^{-1}\!\left(
             r_Q\bigl(\mathcal A(E_Q^{\mathrm c})\bigr)\right).
                                                                  \tag{B.Q8}
\]
For the complete group, with its specified arithmetic subgroup,
\[
 \operatorname{im}D_P=G_{\mathbb Q}\,D_P(H),\qquad
 \operatorname{im}D_P\cap H=D_P(H).                              \tag{B.Q9}
\]
\end{result}

\begin{proof}
First consider $p_Q$.  A full normalized $h\in H$ preserves
$K_Q$ because it is characteristic, and hence preserves its
commutator with $B$.  Its actual pure positive frames reduce
to an element of $\mathcal A(Q^{\mathrm c})$.  Conversely,
the generalized finite-certificate equality above lifts an
element of $\mathcal A(Q^{\mathrm c})$ to a full positive
$\tau$-fixed automorphism $\Phi$ with the indicated $Q$-action.
Write $f=s\Phi(s)$.  The exact normalization (B.M13) changes
$\Phi$ to $\Phi_0=\operatorname{Inn}(x^a)\Phi$, where
$f^{\mathrm{ab}}=(a,-a)$, and supplies an element of $H$.
This changes its $Q$-automorphism only by
$\operatorname{Inn}(q_Q(x)^a)$.  The outer class is therefore
unchanged, proving the first image equality in (B.Q7).

Next take $(\alpha,\beta,A)\in\mathcal A(E_Q^{\mathrm c})$.
Lift its reduction to $E_Q$ to a full positive $\tau$-fixed
automorphism $\Phi$ of $B$, with the prescribed $E_Q$-frames.
For $f=s\Phi(s)$, the literal relations give
\[
\begin{aligned}
 \Phi(x)&=x,\qquad \Phi(y)=f^{-1}yf,\\
 \Phi(z)&=f(z,x)z f(z,x)^{-1},\qquad z=(xy)^{-1}.
\end{aligned}                                                   \tag{B.Q10}
\]
These formulas use the two actual cycle relations and remain
valid before making $f$ a commutator, as proved in B.2.8.
They and the same formulas for $\Phi^{-1}$ show that $\Phi|_F$
preserves the closed normal subgroup
$\ker\varphi_P=\langle\!\langle x^2,z^3\rangle\!\rangle$.
It therefore induces an actual continuous automorphism
$\psi$ of the target $B$.  Since $K_Q$ is characteristic,
$\psi$ preserves $K_Q$.  The square with $q_Q\varphi_P$
then shows that the reduced automorphism of $P_{E_Q}$
preserves $\ker\Theta_Q$ and induces the $Q$-automorphism
obtained by reducing $\psi$.  This proves that $r_Q$ is
defined on its entire displayed domain and is independent
of the chosen full lift.  Reduction, restriction, and
quotient of automorphisms respect composition, so it is a
homomorphism with (B.Q3).

Making $f$ a commutator changes the descended automorphism
by exactly $\operatorname{Inn}(\varphi_P(x)^a)
=\operatorname{Inn}(s^a)$.  Thus the value of $r_Q$ is the
outer $Q$-action of $D_P(h)$ for the resulting $h\in H$.
This proves one inclusion in the second equality of (B.Q7).

For the reverse inclusion, let $h\in H$ be represented by
its normalized $\tau$-fixed automorphism of $B$.  Its
descendant through $\varphi_P$ preserves characteristic
$K_Q$, so $h$ preserves $L_Q$ in (B.Q5).  It also preserves
$B'$ and $F$, and an automorphism takes the $B$-core of a
subgroup to the $B$-core of its image.  Hence $h$ preserves
$K_{E_Q}$ and its closed commutator with $B$.  Its reduction
and its actual pure frames therefore belong to
$\mathcal A(E_Q^{\mathrm c})$, and their $r_Q$-value is
$\rho_Q(D_P(h))$.  This proves the other inclusion.

We justify the passage from all finite levels to (B.Q8).
The family $\{\rho_Q\}$ jointly separates the elements of
$H$.  Indeed equality of two outer actions on every $Q$
means that, in every $Q$, one automorphism is conjugate to
the other by an element of $Q$.  The inverse images in
compact $B$ of these nonempty conjugator sets are closed.
A common characteristic refinement handles every finite
collection of levels.  Compactness gives a single element
of $B$ conjugating the two full automorphisms on $s,r$, and
hence on $B$.  Faithfulness of the normalized representation
$H\to\operatorname{Out}(B)$, proved after (B.M13), identifies
the two gentle elements.

Now let $h$ belong to the right side of (B.Q8).  For each
$Q$, equation (B.Q7) makes the closed set
\[
 \{g\in H:\rho_Q(D_P(g))=\rho_Q(h)\}
\]
nonempty.  These sets have the finite intersection property,
again by passing to a common characteristic refinement.
Compactness of $H$ supplies $g$ belonging to all of them.
The separation just proved yields $D_P(g)=h$.  The reverse
inclusion follows directly from (B.Q7), proving (B.Q8).

Finally $D_P$ preserves the multiplier, while arithmetic
descent (B.C6) gives $D_P(\gamma)=c\gamma c$ for every
$\gamma\in G_{\mathbb Q}$.  Arithmetic cyclotomic
surjectivity gives $\widehat{GT}_0=G_{\mathbb Q}H$ and
therefore the first equality of (B.Q9).  An element of this
image has multiplier one precisely when its displayed
arithmetic factor does.  Such a factor $\gamma$ equals
$D_P(c\gamma c)$ with $c\gamma c\in H$, and hence belongs
to $D_P(H)$.  This proves the second equality.
\end{proof}

\paragraph{The exact frame freedom for a target lift.}
There is useful additional information at the word level.
Let $g\in H$ have its uniquely specified pure positive
frames $\beta,A$, so
\[
 \alpha_g(s)=\beta s\beta^{-1},\quad
 \alpha_g(r)=ArA^{-1},\quad\alpha_g(\tau)=\tau,
 \qquad \beta^{\mathrm{ab}}=(u,u),\quad A^{\mathrm{ab}}=0.
\]
For every actual normalized preimage $(1,f)$ of $g$, its two
evaluation words in the target have the precise form
\[
 \begin{aligned}
 C=f(s,\tau)&=x^{-a}\beta,\\
 J=f(r^{-1},s)&=\beta^{-1}A r^l,
 \qquad a\in\widehat{\mathbb Z},\quad l\in\mathbb Z/3,\\
                 a&\equiv l\equiv -u\pmod3.
 \end{aligned}                                                  \tag{B.Q11}
\]
Here $a$ is the actual target commutator-normalizing exponent
in (B.C13), and $r^l$ is the actual residue of $J$ in $B/F$.
The two displayed word values, for every $a,l$, satisfy the
same raw target automorphism formulas
\[
 s\longmapsto s,\qquad
 \tau\longmapsto C^{-1}\tau C,\qquad
 r\longmapsto JrJ^{-1},
 \qquad \psi=\operatorname{Inn}(\beta^{-1})\alpha_g.             \tag{B.Q12}
\]
Indeed $x$ commutes with $\tau$, and $r^l$ commutes with
$r$.  Thus substitution gives the right side of (B.Q12) on
the generators.  Conversely (B.C13), together with the
pure uniqueness of the two target torsion frames, gives
exactly the two equalities in (B.Q11); no frame is silently
discarded.

The congruences are forced by the actual abelian map
$B\to B^{\mathrm{ab}}=C_6$.  With $[\tau]$ as generator,
the values of $s,r,x,y$ are respectively $3,4,4,4$.
Consequently $\beta$ has value $2u$ and $A$ has value zero.
Since $f$ is a commutator word, each of its two evaluation
words has zero abelian value.  For the words in (B.Q11)
these values are
\[
            -4a+2u=0\pmod6,\qquad
            -2u+4l=0\pmod6,
\]
which are precisely $a\equiv l\equiv-u\pmod3$.
In particular setting both exponents to zero restricts this
construction to targets with $u\equiv0\pmod3$.  It does
not examine the full target set in (B.Q7).

Equations (B.Q7)--(B.Q8) determine the full descent image
from finite marked canonical covers.  The present proof
does not establish equality of the two finite images
$r_Q(\mathcal A(E_Q^{\mathrm c}))$ and
$p_Q(\mathcal A(Q^{\mathrm c}))$ for a cofinal family of
$Q$.  Surjectivity of $D_P$ and independence of all later
descent characters therefore remain unproved here.  The
proved nonzero character in B.2.8 uses its completed
canonical-cover certificate and is unaffected by this
remaining uniform coverage question.
